Our model is built on the framework of RVQ-GANs, following the same pattern as SoundStream\citep{zeghidour2021soundstream}
 and EnCodec\citep{défossez2022high}. As depicted in Figure\ref{fig:model_structure}, our model uses the convolutional-based encoder-decoder network from EnCodec, which performs temporal downscaling with a chosen striding factor. Notably, we have substituted the two-layer LSTM, originally following the convolution blocks in the EnCodec encoder, with a two-layer BiLSTM to augment the semantic modeling ability. We conduct ablation studies of model structure in Appendix~\ref{sec:app:lstm_ablation}. We quantize the encoder outputs using Residual Vector Quantization (RVQ), a method that can operate quantizes residuals following an initial quantization steps with distinct codebook. Further details of model structure can be found in Appendix \ref{sec:app:model}. During training, a semantic teacher provides semantic representation to guide the residual quantization process. %We adopt HuBERT~\citep{hsu2021hubert} as semantic teacher model in this study. 

 
% We represent a single-channel audio signal as a sequence $x \in \mathcal{R}^T$, where $T=d \cdot f_{sr}$ denotes the length with duration $d$ at sampling rate $f_{sr}$.
% As illustrated in Figure~\ref{fig:model_structure}, \method model consists of four main components: 
% \begin{itemize}
%     \item an encoder network $E$ that takes an audio as input and output a latent representation $z$.
%     \item a residual vector quantizer $Q$ that compresses representation $z$ to $z_q$ from a set of finite codebooks.
%     \item a decoder network $G$ that produces the reconstructed audio signal $\hat{x}$ from the compressed representation $z_q$.
%     \item a semantic teacher model $S$ that provides semantic representation $s$ to guide the residual vector quantization process during training.
% \end{itemize}
% \method is trained end-to-end with reconstruction loss, together with a perceptual loss in the form of discriminators and semantic distillation loss as training objective. We adopt HuBERT~\citep{hsu2021hubert} as semantic teacher model in this study. 

% \textbf{Encoder \& Decoder Architecture}~
% Inspired by SoundStream~\citep{zeghidour2021soundstream} and EnCodec~\citep{défossez2022high}, our model employs a convolutional-based encoder-decoder architecture. The encoder is constructed as a sequential series of components: starting with a 1D convolutional layer featuring $C$ channels and a kernel size of 7, followed by a set of $B$ residual conventional blocks. Each block is composed of two dilated convolutions with dilation rate of $(1,1)$ and kernel size of $(3,1)$ and a skip-connection, followed by a strided convolutional down-sampling layer, with a kernel size $K$ of the twice the stride $R$. Whenever down-sampling, the number of channels is doubled. Unlike in EnCodec that the convolution blocks are followed by a two-layer LSTM, we use BiLSTM to augment the semantic modeling ability. A final 1D convolution layer with a kernel size of 7 is used to set the dimensionality of embeddings to $D$. We use $C=32,B=4$ and $(2,4,5,8)$ as strides. We use ELU~\citep{clevert2016fast} as a non-linear activation either layer normalization~\citep{ba2016layer} or weight normalization~\citep{salimans2016weight}.The decoder mirrors the encoder and uses transposed convolutions and LSTM instead of stride convolutions and BiLSTM, with the strides in reverse order as in the encoder. The decoder outputs the final audio signal.

% \textbf{Residual Vector Quantizer}~
% We use Residual Vector Quantizer~(RVQ) to quantize the encoder output. Vector quantization~(VQ) involves selecting the code closest to the input vector from a codebook of a given size as the quantization representation. Residual vector quantizer cascades $N_q$ layers of VQ to alleviate the conflict between codebook size and quantization performance. The unquantized input vector is passed through the first quantizer and quantization residuals are computed. These residuals are then iteratively quantized by a sequence of $N_q-1$ additional quantizers. For each residual step $i\in\{1,\cdots, N_q\}$, $z_i$ denotes the quantize-before vector and $q_i$ denotes quantize-after vector and $c_i$ denotes the corresponding code index.

% We follow the same training procedure as EnCodec. The code selected for each input is updated using an exponential moving average with a decay of 0.99, and codes which have not been assigned any input vector for several batches are replaced with input vectors randomly sampled within current batch. Straight-through-estimator~\citep{bengio2013estimating} is used to compute the gradient of encoder, e.g. as if the quantization step was the identity function during the backward phase. Finally, a commitment loss, consisting of the MSE between the input of the quantizer and its output, with gradient only computed with respect to its input, is added to the overall training loss.

% \textbf{Discriminator}~
% We use the same dicriminators as HiFi-Codec~\cite{yang2023hificodec} that consist of three discriminators: A multi-scale STFT-based~(MS-STFT) discriminator; a multi-period discriminator~(MPD) and a multi-scale discrinator~(MSD). The MS-STFT discriminator utilizes networks with identical structures that operate on multi-scaled complex-valued STFT, where the real and imaginary parts are concatenated. For each sub-network, it is composed of a 2D convolutional layer (using kernel size $3\times8$ with 32 channels), followed by 2D convolutions with increasing dilation rates in the time dimension (1, 2 and 4), and a stride of 2 over the frequency axis. A final 2D convolution with kernel size $3\times3$ and stride $(1, 1)$ provide the final prediction. For MSD and MPD, we follow the same settings as in HiFiGAN~\citep{kong2020hifi} but adjust the channel number to align the discriminator's parameters more closely with that of MS-STFT.